% 这是一份全新的个人教程副本；修改它不会影响任何其他人的项目。
% 主文件：main.tex；编译引擎：XeLaTeX。
\documentclass[UTF8,fontset=fandol,11pt,a4paper]{ctexart}
\usepackage[margin=24mm]{geometry}
\usepackage{amsmath,amssymb,graphicx,booktabs,xcolor,hyperref,fancyhdr}
\definecolor{Ink}{HTML}{283D3B}
\definecolor{Light}{HTML}{F4F1E7}
\hypersetup{colorlinks=true,linkcolor=Ink,urlcolor=Ink,citecolor=Ink,pdftitle={纸上微光 · Prims 写作手册}}
\setlength{\parindent}{2em}
\setlength{\parskip}{0.4em}
\setlength{\headheight}{15pt}
\pagestyle{fancy}
\fancyhf{}\fancyhead[L]{\small\color{Ink}纸上微光}\fancyhead[R]{\small Prims}\fancyfoot[C]{\thepage}
\newcommand{\menu}[1]{\textbf{#1}}
\newcommand{\tryit}[1]{\par\smallskip\noindent\colorbox{Light}{\parbox{\dimexpr\linewidth-2\fboxsep\relax}{\textbf{动手试试}\quad #1}}\par\smallskip}
\title{\vspace{-1em}{\Huge\color{Ink}纸上微光}\\[0.5em]{\Large Prims 写作手册}}
\author{从一页纸，到一份作品}
\date{2026 年 10 月版}
\begin{document}
\maketitle
\begin{center}\begin{minipage}{0.84\linewidth}
写作常常从一个还不清楚的念头开始。你可以先写下来，再和助手讨论；先改一句，再慢慢整理整篇。

这份手册也是一个可以编译、编辑和撤销的项目。它属于你，旧项目和旧教程都保留原样。
\end{minipage}\end{center}
\tableofcontents
\section{先写下一页}
\subsection{把作品带进来}
项目首页左侧是书架，选中一本，右侧可以翻阅最近编译的 PDF；点击\menu{进入项目}开始编辑。书架可以独立滚动，项目旁的菜单用于改名或删除。

\menu{上传／导入}可以带入单个文件、多个文件、文件夹或 ZIP 压缩包。导入文件夹时保留内部目录，ZIP 会展开为项目文件。已有项目里也可以继续上传。文件夹选择需要浏览器支持，空目录不会随文件夹导入；ZIP 请单独选择。遇到同名路径时先改名再导入，已有文件不会被静默覆盖。

左侧文件列表支持新建文件、目录、改名和拖动移动。开启\menu{批量选择}后可以导出、移动或删除多个文件；处理结果会逐项显示。删除图片等附件前，先导出一份副本，文本历史不能代替二进制附件备份。

\subsection{源文件与 PDF}
编辑区里的 \texttt{main.tex} 是这份手册的主文件，正文通过 \verb|\input| 引入 \texttt{chapters/} 下的各章。点击\menu{编译}后，右侧 PDF 展示整个项目的结果；单独打开某一章，不会把编译入口自动换成那一章。

需要换入口时，在项目设置中选择\menu{主文件}。同一处可以选择编译引擎：中文或自定义字体项目通常用 XeLaTeX；已有模板则先沿用作者的设置。这个手册已经选择了 XeLaTeX，直接编译即可。

\tryit{打开本章文件，把下面一句改成你想写的题目，等自动保存完成后点击编译。}
\begin{quote}我正在写一份关于人与城市的文档。\end{quote}

\subsection{把阅读环境调舒服}
\menu{设置 → 外观}里可以选一整套日夜配色，切换浅色、深色或跟随系统，也可以调整界面字体、书名字体和 Prims 字标。它们不会改变论文的 LaTeX 字体或排版。

PDF 区域的外观下拉可以选择阅读效果，包括暖黄或浅绿护眼色。阅读滤镜只影响屏幕显示，下载出来的 PDF 仍保留文档原本的颜色。

\section{让助手接上你的思路}
\subsection{先接好模型}
打开\menu{设置 → 个人 AI 配置}，新建一个命名连接，填写接口地址、API Key 和模型名称。一个连接可以填写多个模型，每行一个；不同服务商可以分别建连接。接口格式能识别时自动选择，不能识别时，按服务商文档选择 OpenAI 或 Anthropic 等对应格式。

提问框里的模型菜单可以切换本轮使用的模型。密钥属于你自己的账号，不与其他成员共享。各模型的上下文长度也可以单独填写：窗口越大，能携带的上下文越多，同时可能增加单次请求用量。若接口提示上下文过长，对话中的卡片可以直接调整；用量图在个人 AI 配置里按当天、七天或一个月查看。

\subsection{讨论与修改都可以}
直接在助手里说你正在做什么，例如：\emph{我在写实验设计，想让这一段更清楚地说明对照条件。}不必先想好一条完美提示词。只是问问题时，助手可以给文字解释，不一定产生修改。

选中文字后会出现一个小输入框。选区是定位线索，助手还可以读取当前文件、按需查找项目里的其他文件与图片，并提出跨文件建议。把关键要求说清楚，比反复粘贴整篇内容更有效。

\tryit{选中上一节的示例句，试着问：“保留意思，把这句话写得具体一点；如果信息不够，先问我。”}

建议会显示红色删除与绿色新增。你可以先改候选文字，再逐段\menu{接受}或\menu{拒绝}；需要时也可以全部接受。接受前正文不会改变，已经接受的修改可以撤销。若你在等待时又改过同一段，先核对冲突提示，再决定采用哪份文字。

\subsection{从一句话走到整个项目}
助手能按需列目录、搜索内容、读文件，提出新建文件或目录、移动文件、修改编译设置等项目操作。比如：\emph{把引言独立为一个章节，并补上主文件中的引用。}先让它看真实文件，再检查提出的文字与文件操作，接受后才会写入。

图片需要模型具备视觉能力。如果当前模型不支持，系统会尝试去掉图片继续；要讨论图中的细节，请切换到能看图的模型。

\menu{设置 → Agent}管理搜索连接、技能和项目记忆。搜索的接口类型、endpoint 与 key 按搜索服务商填写；助手会在需要时调用。技能可以整理成反复使用的写作流程，也可以从网页导入。技能里的文字不会因此获得额外的系统权限。项目记忆适合放研究目标、术语约定与长期偏好，不必把整篇论文重复写进去。

\tryit{试着说：“先梳理本项目的文件结构，再告诉我哪些章节适合补一个例子；暂时只讨论。”这样可以练习让助手先了解作品，再进入修改。}

\section{在 LaTeX 工作台上}
\subsection{标题、公式与交叉引用}
正文只需专注内容，章节结构交给 \verb|\section| 和 \verb|\subsection|。公式写在数学环境中；例如，式~\ref{eq:mean} 表示一组观测的算术平均数：
\begin{equation}
  \bar{x}=\frac{1}{n}\sum_{i=1}^{n}x_i .
  \label{eq:mean}
\end{equation}
用 \verb|\label| 给公式或章节起一个稳定标记，再用 \verb|\ref| 引用它。改动顺序后，重新编译会更新编号。表~\ref{tab:files} 是这份项目的一个小目录。
\begin{table}[htbp]
\centering
\caption{手册里的几个文件}\label{tab:files}
\begin{tabular}{ll}\toprule
文件 & 用途\\\midrule
\texttt{main.tex} & 入口、样式与章节顺序\\
\texttt{chapters/} & 分章正文\\
\texttt{refs.bib} & 参考文献条目\\\bottomrule
\end{tabular}
\end{table}

\subsection{参考文献与图片}
本项目的 \texttt{refs.bib} 中放了一条真实的 LaTeX 官方文档入口，可以用 \verb|\cite{latex-docs}| 引用\cite{latex-docs}。新文献可以由你或助手创建条目，但作者、题名、年份和链接仍应核对来源，尤其不要把示例条目当成自己的研究证据。

把图片上传到项目后，可在正文里使用 \verb|\includegraphics{figures/example.png}|。路径要与文件列表一致；图片文件与正文一起保存在项目中。需要 AI 新建 \texttt{refs.bib}、整理图目录或调整引用时，直接说明目标，再检查它的文件操作建议即可。

\subsection{字体与编译引擎}
这份中文手册用 \texttt{ctexart} 与 TeX 自带的 Fandol 字体，尽量让第一次编译就顺利。项目的文档字体由 LaTeX 源码控制，和网页设置里的编辑器字体是两回事。

使用 XeLaTeX 或 LuaLaTeX 时，可以通过 \texttt{fontspec} 指定编译环境可用的字体。可以请助手先查询可用字体，再根据返回的字体名修改设置。若另一个模板提示找不到字体，先核对字体名；需要随项目携带字体时，可以上传你有权使用的 \texttt{.ttf} 或 \texttt{.otf} 文件，再按项目路径引用。不必为一个字体错误重写正文。

\subsection{出错时，从第一条错误开始}
编译失败后切到\menu{日志}，先看第一条有效错误及其文件、行号；后面一长串往往只是连带结果。检查括号、环境配对、文件路径和字体名，通常比从头重写更快。

可以把编译错误交给助手处理。它会结合日志与项目文件做诊断、提出修改，也能建议更换主文件或引擎。逐项确认后再编译；若修改不合适，可以拒绝候选，或从版本历史找回先前文本。

\tryit{先在本章末尾加一个练习段，试写公式、表格或一段引文。保持主文件为 \texttt{main.tex}，重新编译后查看效果。}

\section{留住每一次向前}
\subsection{自动保存与版本历史}
正常编辑会自动保存；离开前留意保存状态，网络断开时先保留本地文字。\menu{版本历史}按日期和时间展示项目的文本版本，可以查看差异、恢复内容，也可以让 AI 为改动生成摘要。摘要基于改动内容，不需要把整份作品重新讲一遍。

\menu{撤销}适合回退刚刚发生的一步，版本历史适合回到更早的时间点。两者都以实际保存的范围为准；二进制图片和附件建议另外导出留存。恢复前先看预览，避免把后来写好的内容一起换掉。

\subsection{把助手放在顺手的位置}
助手可以停靠在编辑器右侧、底部，也可以浮动。位置与大小可以按习惯调整。需要更宽的正文时，把助手放到独立标签页；它仍会跟随编辑器当前打开的文件和选区，也可以再放回编辑页。

选区小输入框与固定助手的输入草稿互不打扰，发送后仍使用同一段对话与上下文。对话标题处可以展开历史列表，新建对话不会删除旧讨论；需要接着写时，再选回原来的对话。

\subsection{一个轻一点的写作节奏}
先用一两句话说明你正在完成什么，给助手读真实材料的机会；让它先做一小段，再把有效的改法推广到相关位置。每次接受后看一眼正文和 PDF，需要停下来思考时，就只讨论。

这份手册可以留在书架上，也可以改名、删掉练习段落，变成你自己的第一份作品。工具尽量安静，写作仍在你手里。

\bibliographystyle{plain}
\bibliography{refs}
\end{document}
